\subsection{问题三的建模与求解：逐次往返厚度对照}

上一问给出了碳化硅的两束模型基准，本问采用逐次往返厚度对照，在相同设置下比较两束截断场与完整往返场。完整往返场将表面反射与各轮返回相加，往返衰减场反演据此计算反射率，再搜索与观测相符的厚度。硅在这一框架下求厚度，碳化硅则在高阶往返可观测且影响厚度时修正。

完整往返场给出硅全量厚度 $2.17\ \mu\mathrm{m}$，两次连续留段分别为 $1.73$、$16.04\ \mu\mathrm{m}$；碳化硅缺少相干、空间重叠和仪器分辨率资料，保留前问基准 $10.80\ \mu\mathrm{m}$。
% src:求解/问题3/结果/回炉轮3候选/仲裁闭环/20260910_131501_81623_建模公式/厚度结果.json:核心指标.硅全量完整往返厚度_um;核心指标.硅折1完整往返条件厚度_um;核心指标.硅折2完整往返条件厚度_um;核心指标.碳化硅正式基准厚度_um

两材料在同波段、同留出点及同响应设置下比较，完整往返场的总体标准化留段均方根误差为 $1.0231$，高于两束的 $0.9812$（均无量纲）。逐块收益随材料和区段改变，完整场并未取得总体预测优势。这个分歧要求先核对观测条件与资料缺口，再推导往返场及其厚度影响。
% src:求解/问题3/结果/回炉轮3候选/仲裁闭环/20260910_131501_81623_建模公式/厚度结果.json:分材料验证.硅.逐折汇总.0.平均标准化均方根误差.两束;分材料验证.硅.逐折汇总.0.平均标准化均方根误差.完整往返;分材料验证.硅.逐折汇总.0.平均标准化均方根误差.训练均值;分材料验证.硅.逐折汇总.0.平均标准化均方根误差.二次趋势;分材料验证.硅.逐折汇总.1.平均标准化均方根误差.两束;分材料验证.硅.逐折汇总.1.平均标准化均方根误差.完整往返;分材料验证.硅.逐折汇总.1.平均标准化均方根误差.训练均值;分材料验证.硅.逐折汇总.1.平均标准化均方根误差.二次趋势;分材料验证.碳化硅.逐折汇总.0.平均标准化均方根误差.两束;分材料验证.碳化硅.逐折汇总.0.平均标准化均方根误差.完整往返;分材料验证.碳化硅.逐折汇总.0.平均标准化均方根误差.训练均值;分材料验证.碳化硅.逐折汇总.0.平均标准化均方根误差.二次趋势;分材料验证.碳化硅.逐折汇总.1.平均标准化均方根误差.两束;分材料验证.碳化硅.逐折汇总.1.平均标准化均方根误差.完整往返;分材料验证.碳化硅.逐折汇总.1.平均标准化均方根误差.训练均值;分材料验证.碳化硅.逐折汇总.1.平均标准化均方根误差.二次趋势
% src:交接/实验记录.json:563.尝试;563.现象;563.决定
